% VI/36 reader bridge -- additive, no replacement of canonical material.
\paragraph{From an affine chart to projective space.}
Let \(k\) be a field.  A \(k\)-rational point of projective space is
represented by a nonzero tuple
\((a_0,\ldots,a_n)\in k^{n+1}\), where
\((a_0,\ldots,a_n)\) and \((\lambda a_0,\ldots,\lambda a_n)\)
represent the same point for every \(\lambda\in k^\times\).
We write its equivalence class as \([a_0:\cdots:a_n]\).
This rule is a useful description of \(\mathbb P^n(k)\), but the
scheme \(\mathbb P^n_k=\operatorname{Proj}k[X_0,\ldots,X_n]\)
has additional points with nontrivial residue-field extensions.
The distinction prevents arguments on coordinate tuples alone from
being mistaken for statements about all scheme-theoretic points.

\paragraph{Worked example: three affine charts of \(\mathbb P^2\).}
The homogeneous polynomial ring is
\(S=k[X,Y,Z]\) with each variable of degree one.
On \(U_Z=D_+(Z)\), the coordinate functions are the ratios
\(u=X/Z\) and \(v=Y/Z\), so
\[
 U_Z\simeq\operatorname{Spec}(S_Z)_0
       =\operatorname{Spec}k[u,v]\simeq\mathbb A_k^2.
\]
Likewise \(U_X\simeq\operatorname{Spec}k[Y/X,Z/X]\) and
\(U_Y\simeq\operatorname{Spec}k[X/Y,Z/Y]\).
The three opens cover the scheme, since a homogeneous prime in
\(\operatorname{Proj}S\) cannot contain all of \(X,Y,Z\).
Where both \(X\) and \(Z\) are nonzero, the coordinates change by
\[
  \frac{Y}{X}=\frac{Y/Z}{X/Z}=\frac{v}{u},\qquad
  \frac{Z}{X}=\frac{1}{u}.
\]
These reciprocal expressions are regular on the \emph{overlap},
where \(u\) is invertible; they are not regular functions on all of
\(\mathbb A_k^2\).  This is the concrete gluing mechanism for
projective space.

\paragraph{A parabola and the point at infinity.}
Inside the chart \(Z\ne0\), the affine parabola has equation
\(v=u^2\).  Substituting \(u=X/Z\) and \(v=Y/Z\) and clearing the
homogeneous denominator of degree two gives
\[
 YZ-X^2=0.
\]
Its projective closure is the closed subscheme
\(C=V_+(YZ-X^2)\subset\mathbb P_k^2\).
To inspect the complement of the affine chart, put \(Z=0\).
Then \(X^2=0\); at the level of points \(X=0\), and the only
projective point on the hyperplane \(Z=0\) is \([0:1:0]\).
The scheme-theoretic intersection with the hyperplane is defined by
\(Z=X^2=0\), and may have nonreduced structure: on the chart \(Y=1\)
its ring is \(k[x]/(x^2)\).  Thus there is one point at infinity,
but the intersection multiplicity is two.
This is why point counting alone is insufficient for projective
intersection problems.

\paragraph{A morphism to projective space must be defined everywhere.}
The quadratic formula
\[
 \nu_2([s:t])=[s^2:st:t^2]
\]
uses homogeneous forms of the same degree, and they cannot all
vanish at a projective point.  It therefore defines a morphism
\(\mathbb P_k^1\to\mathbb P_k^2\), not merely a rational map.
On \(s\ne0\), set \(w=t/s\); then
\(\nu_2([1:w])=[1:w:w^2]\).
The image satisfies \(XZ=Y^2\).
VI/38 verifies the stronger conclusion that the morphism induces
an isomorphism onto this \emph{closed subscheme}; checking only the
homogeneous equation and pointwise injectivity would not suffice.

\paragraph{Why the Segre coordinate rule is well defined.}
For \(\mathbb P_k^1\times_k\mathbb P_k^1\), consider
\[
 ([a:b],[c:d])\longmapsto[ac:ad:bc:bd]\in\mathbb P_k^3.
\]
Rescaling the first tuple by \(\lambda\) and the second by
\(\mu\) rescales all four products by the single factor
\(\lambda\mu\).  Therefore the class is independent of the chosen
representatives.  Its image satisfies the homogeneous relation
\(Z_{00}Z_{11}-Z_{01}Z_{10}=0\).
The corresponding affine-chart argument that makes this a closed
embedding, including its local ring maps, belongs to VI/38.
